% 多分类维度关注对象的层次：标签对、参考标签、局部邻接。
\begin{tikzpicture}[
  panel/.style={book node,minimum width=31mm,minimum height=42mm,inner sep=4pt},
  tag/.style={draw=BookRule,rounded corners=1pt,minimum width=6mm,minimum height=5mm,font=\scriptsize},
  note/.style={font=\kaishu\scriptsize,text=BookMuted,align=center}
]
  \node[panel,fill=FigureInput!7] (nat) {};
  \node[panel,fill=FigureModel!7,right=6mm of nat] (graph) {};
  \node[panel,fill=FigureLoss!7,right=6mm of graph] (ds) {};

  \node[font=\heiti\small,anchor=north] at ([yshift=-3mm]nat.north) {Natarajan维};
  \node[font=\heiti\small,anchor=north] at ([yshift=-3mm]graph.north) {图维};
  \node[font=\heiti\small,anchor=north] at ([yshift=-3mm]ds.north) {DS维};

  \foreach \i/\x in {1/-9,2/0,3/9}{
    \node[tag,fill=white] at ([xshift=\x mm,yshift=6mm]nat.center) {$a_{\i}$};
    \node[tag,fill=white] at ([xshift=\x mm,yshift=-4mm]nat.center) {$b_{\i}$};
    \draw[BookBlue,thick,<->] ([xshift=\x mm,yshift=3mm]nat.center) -- ([xshift=\x mm,yshift=-1mm]nat.center);
  }
  \node[note,anchor=south] at ([yshift=2mm]nat.south) {每个坐标预先固定\\一对可切换标签};

  \foreach \i/\x in {1/-9,2/0,3/9}{
    \node[tag,fill=white] at ([xshift=\x mm,yshift=6mm]graph.center) {$r_{\i}$};
    \node[tag,fill=white] at ([xshift=\x mm,yshift=-4mm]graph.center) {$\neq r_{\i}$};
    \draw[FigureModel,thick,->] ([xshift=\x mm,yshift=3mm]graph.center) -- ([xshift=\x mm,yshift=-1mm]graph.center);
  }
  \node[note,anchor=south] at ([yshift=2mm]graph.south) {固定参考标签；离开后\\具体标签可以变化};

  \node[tag,fill=white] (v1) at ([xshift=-8mm,yshift=5mm]ds.center) {$abc$};
  \node[tag,fill=white] (v2) at ([xshift=8mm,yshift=5mm]ds.center) {$dbc$};
  \node[tag,fill=white] (v3) at ([xshift=-8mm,yshift=-6mm]ds.center) {$aec$};
  \node[tag,fill=white] (v4) at ([xshift=8mm,yshift=-6mm]ds.center) {$abf$};
  \draw[FigureLoss,thick] (v1)--(v2);
  \draw[FigureLoss,thick] (v1)--(v3);
  \draw[FigureLoss,thick] (v1)--(v4);
  \node[note,anchor=south] at ([yshift=2mm]ds.south) {研究完整标记向量间\\只改一个坐标的邻接};
\end{tikzpicture}
